\section{Formulación del problema}\label{sec:formulacion}

El problema es una decisión estática: elegir en $t_0$ (31-12-2025) el monto de cada posición de activo y de
pasivo para el horizonte de un año ($t_H$ = 31-12-2026), sin rebalanceo dentro del año. Se optimiza la
estructura conjunta de los seis activos y los seis pasivos existentes; no se considera deuda nueva, porque el
caso asignado no la incluye en su alcance. El Anexo~\ref{app:formulacion} desarrolla la formulación completa y los supuestos en que se apoya.

\subsection{Variables de decisión}
Sea $i$ un activo, $j$ un pasivo y $k$ cualquiera de los doce instrumentos; $z_k$ denota indistintamente
$x_i$ o $y_j$ y $z_k^0$ la posición inicial. Las variables son:

\begin{center}
\begin{tabular}{clll}
\toprule
Variable & Descripción & Dominio & Unidad \\
\midrule
$x_i$ & monto asignado al activo $i$ & $x_i \ge 0$ & S/ mm \\
$y_j$ & monto asignado al pasivo $j$ & $y_j \ge 0$ & S/ mm \\
$\Delta_k^+,\ \Delta_k^-$ & aumento y reducción frente a la posición inicial $z_k^0$ & $\ge 0$ & S/ mm \\
$\zeta$ & valor en riesgo auxiliar (Rockafellar--Uryasev) & libre & S/ mm \\
$u_s$ & exceso de pérdida del escenario $s$ sobre $\zeta$ & $\ge 0$ & S/ mm \\
\bottomrule
\end{tabular}
\end{center}

No hay posiciones cortas: la ausencia de esa alternativa en los datos del caso se traduce en $x_i, y_j \ge 0$.
La reestructuración de cada instrumento se descompone en $z_k - z_k^0 = \Delta_k^+ - \Delta_k^-$, lo que
permite penalizar simétricamente cualquier movimiento, subida o bajada, sin usar variables binarias.

\subsection{Función objetivo}
El resultado del patrimonio y su pérdida en el escenario $s$ son
\begin{equation}
  \Delta \PN_s = \sum_{i} x_i\, r_{i,s} - \sum_{j} y_j\, r_{j,s}, \qquad L_s = -\Delta \PN_s,
\end{equation}
donde $r_{k,s}$ es el resultado a un año del instrumento $k$ bajo el escenario $s$ (Sección~\ref{sec:riesgo}). El
costo de transacción $TC = \sum_k c_k\,(\Delta_k^+ + \Delta_k^-)$, donde $c_k$ es el costo de
reestructuración por sol movido del instrumento $k$ que fija el caso, se aplica igual a subidas y bajadas.
El costo de prepago no se suma a las reducciones de pasivos, porque el caso lo reserva para otros
escenarios de decisión. $TC$ se paga en $t_0$, es determinista y no depende del escenario. El problema es

\begin{equation}
  \max_{x,\,y,\,\Delta,\,\zeta,\,u}\quad \sum_{s} p_s\,\Delta \PN_s \;-\; \lambda\left(\zeta + \frac{1}{1-\alpha}\sum_s p_s\, u_s\right) \;-\; TC
  \qquad\text{s.a. } u_s \ge L_s - \zeta,\ u_s \ge 0.
  \label{eq:objetivo}
\end{equation}

El término entre paréntesis es el $\CVaR_\alpha$ de la pérdida $L_s$ en su forma lineal de
\textcite{rockafellar2000}: minimizarlo sobre $\zeta$ recupera exactamente el valor en riesgo al nivel
$\alpha$ y el promedio de las pérdidas que lo superan. Como el objetivo lo resta, ponderado por la
aversión al riesgo $\lambda \ge 0$, al maximizar el óptimo elige el $\zeta$ que minimiza esa expresión, y la
cola queda penalizada sin perder la linealidad del problema. Con $\lambda = 0$ el objetivo maximiza solo el resultado esperado neto de costos; a medida que
$\lambda$ crece, la solución se acerca a la que minimiza el $\CVaR_{0.95}$ alcanzable (Sección~\ref{sec:resultados}).
El costo de transacción se resta una sola vez, fuera de $L_s$: es un costo cierto de la decisión, no un
riesgo de mercado, de modo que se reporta por separado el patrimonio neto de costos, $\PN_0' = \PN_0 - TC$,
y su valor esperado al horizonte, $\E[\PN_H] = \PN_0' + \E[\Delta \PN]$. La rentabilidad esperada de los
activos es $\sum_s p_s \sum_i x_i r_{i,s} / B_A$ y el costo esperado de los pasivos,
$\sum_s p_s \sum_j y_j r_{j,s} / B_L$. Este es un costo económico: incluye intereses y revalorización.

Dos variantes sirven de comparación. La cartera de \emph{mínima varianza} minimiza
$\sum_s p_s (\Delta \PN_s - \E[\Delta \PN])^2$ sujeta a $\E[\Delta \PN] - TC$ mayor o igual que el de la
cartera óptima y a las mismas restricciones. La variante con \emph{vencimientos exigidos al horizonte}
agrega, para cada grupo de vencimiento remanente, la misma cota de la ecuación~\eqref{eq:pesos} medida con
el plazo que le queda a cada instrumento en $t_H$.

Con probabilidades no uniformes y masas de probabilidad concentradas en unos pocos escenarios de estrés, el
cuantil que define el valor en riesgo puede no ser único cuando $\lambda \to 0$; por eso el $\VaR$ y el
$\CVaR$ que se reportan en los resultados se recalculan \emph{ex post}, sobre la distribución completa de
$\Delta \PN_s$, con la definición de átomo fraccional
$\CVaR_\alpha = \frac{1}{1-\alpha}\big[\sum_{L_s > \VaR} p_s L_s + \VaR\,(1-\alpha - \sum_{L_s > \VaR} p_s)\big]$,
en vez de leer directamente la variable $\zeta$ de la ecuación~\eqref{eq:objetivo}.

\subsection{Restricciones}
Los presupuestos $\sum_i x_i = B_A = 1{,}000$ y $\sum_j y_j = B_L = 800$ (S/ mm) son exógenos y fijan el
denominador de cada peso, de modo que las restricciones del Cuadro~4 del enunciado son lineales en las variables
de decisión: para cada grupo $g$ (moneda, tipo de instrumento, tipo de tasa o vencimiento remanente) y cada
lado $l \in \{A, L\}$,
\begin{equation}
  \underline{w}_g\, B_l \;\le\; \sum_{k \in \mathcal{K}_g} z_k \;\le\; \overline{w}_g\, B_l.
  \label{eq:pesos}
\end{equation}
Un grupo adicional, de concentración, acota individualmente cada instrumento ($z_k \le \overline{w}_g\, B_l$).
Sobre esos límites, el patrimonio neto en $t_0$ debe ser de al menos \PEN120 mm, el cociente pasivos/activos
de a lo más 0.85 y la caja de al menos \PEN80 mm; con los presupuestos fijos, las dos primeras ya se cumplen
por construcción ($\PN_0 = 200$, $P/A = 0.80$) y se mantienen en el modelo solo para que el solver las
verifique y el informe las reporte.

\subsection{Modelo de riesgo}\label{sec:riesgo}
El resultado por instrumento $r_{k,s}$ no se aproxima con una duración o un DV01, sino que se obtiene
revalorizando cada instrumento completo al horizonte bajo el escenario $s$:
\begin{equation}
  r_{k,s} = \frac{V_k^{PEN}(t_H;\,\text{mercado}_s) + CF_k^{PEN}(t_0, t_H]_s}{V_k^{PEN}(t_0)} - 1,
\end{equation}
donde $V_k(t_H)$ es el valor presente de los flujos posteriores a $t_H$ descontados con la curva del
escenario y $CF_k(t_0, t_H]$ son los flujos cobrados o pagados dentro del año, convertidos a soles al tipo de
cambio del escenario en su fecha de pago y sin reinvertir. Como $r_{k,s}$ no depende del tamaño de la
posición, la ecuación~\eqref{eq:objetivo} con esta matriz de resultados fijada de antemano es un problema
lineal.

Los escenarios $s$ combinan dos fuentes, con probabilidad conjunta 95\,\%/5\,\%: 500 trayectorias de un año
obtenidas por remuestreo por bloques de la historia mensual \parencite{kunsch1989}, y los cuatro escenarios de estrés propios del caso, cada uno
con probabilidad 1.25\,\%. El remuestreo concatena dos bloques de seis meses consecutivos de cambios
históricos, en vez de doce meses independientes, para conservar la persistencia mensual documentada en la
Sección~\ref{sec:datos} y no subestimar la varianza anual. A ese shock centrado se le suma una deriva anual que
depende de la clase de factor: nula en las tasas (para no proyectar el ciclo alcista de 2021--2025 hacia
adelante), la paridad de tasas implícita en las curvas de $t_0$ en el tipo de cambio, y el promedio histórico
del log-retorno mensual en la renta variable (aproximadamente 6\,\%/año). En el tipo de cambio, el shock
se desplaza además por una constante para que la media de $FX_H/FX_0$ en la muestra sea exactamente la
forward de paridad. Sumar la paridad al logaritmo sin este ajuste la sobreestima por convexidad
($\E[e^X] = e^{\mu + \sigma^2/2}$) y por el error de muestreo de la media. Dentro del horizonte, cada factor se
mueve linealmente entre su valor en $t_0$ y su valor bajo el escenario en $t_H$, la misma regla con la que se
reproyectan los cupones flotantes que se fijan dentro del año y con la que devenga la caja.
